\section{Robustness to Prior Distribution Shape}\label{sec:exp7-prior-shape}

The prior misspecification experiment of Section~\ref{sec:exp1-misspec}
tested what happens when the prior probabilities $p_i$ seen by the auction
are noisy corruptions of the true prior from which the target was drawn.
A complementary and arguably more fundamental robustness question is
whether the \emph{model ranking}---M2 $<$ M4* $\leq$ M4 $<$ M3 $<$
H$_d$ $<$ M1---holds when the prior distribution itself has a
qualitatively different shape, not merely when it is noisy.
In a real deployment the analyst may not know which prior shape is
most appropriate; the robustness of the ranking across shapes therefore
determines whether the conclusions of this thesis transfer to
deployments with different prior models.

Five prior shapes were tested over 500 paired trials each at the default
configuration ($n = 30$, $R = 3$, $E = 14$, $L = 10$):

\begin{enumerate}
    \item \emph{Concentrated} ($\mathrm{Dirichlet}(\alpha = 0.2)$): one
          or two sites dominate, modeling a scenario where strong evidence
          concentrates the prior.
    \item \emph{Moderate} ($\mathrm{Uniform}(0.1, 1.0)$, normalized):
          the paper's baseline---mild variation, no strong peak.
    \item \emph{Diffuse} ($\mathrm{Dirichlet}(\alpha = 2.0)$): broader
          spread with no dominant site.
    \item \emph{Near-uniform} ($\mathrm{Dirichlet}(\alpha = 10.0)$):
          essentially uniform, with very weak prior information.
    \item \emph{Power-law} ($U(0,1)^3$, normalized): heavy-tailed,
          modeling sparse but strong indications at a few sites.
\end{enumerate}

The target was resampled from each prior subject to the same geometric
feasibility filters ($d_{\mathrm{opt}} \geq 1.0$,
$2\,d_{\mathrm{opt}} \leq E$) used throughout the thesis.

\begin{figure}[ht]
  \centering
  \includegraphics[width=1.0\textwidth]{fig_prior_sensitivity.png}
  \caption{Mean finder competitive ratio by model and prior distribution
    shape, over 500 paired trials per shape at the default configuration
    ($n = 30$, $R = 3$, $E = 14$, $L = 10$).
    Three of five prior shapes preserve the complete ordering
    M2 $<$ M4* $\leq$ M4 $<$ M3 $<$ H$_d$ $<$ M1.
    The concentrated prior reveals a regime where M4* surpasses M2:
    the quadratic penalty guides the robot directly to the single dominant
    high-probability site so efficiently that the energy constraint
    becomes irrelevant.
    The power-law prior reveals a narrow exception where M4 outperforms
    M4*: extreme skew in the prior creates a few very high-probability
    distant sites that the quadratic penalty discounts too
    aggressively.}\label{fig:prior-sensitivity}
\end{figure}

\begin{table}[ht]
\centering
\caption{Mean finder competitive ratio by prior distribution shape and
  model (500 paired trials per row, $n = 30$, $R = 3$, $E = 14$).
  A checkmark indicates the full ordering
  M2 $<$ M4* $\leq$ M4 $<$ M3 $<$ H$_d$ $<$ M1 holds;
  a cross indicates at least one pair is
  violated.}\label{tab:prior-sensitivity}
\vspace{2pt}
\begin{tabular}{@{}lrrrrrrc@{}}
\toprule
\textbf{Prior Shape} & \textbf{M2} & \textbf{M4*} & \textbf{M4}
  & \textbf{M3} & \textbf{H$_d$} & \textbf{M1} & \textbf{Ordering} \\
\midrule
Concentrated (Dir~0.2) & 1.89 & \textbf{1.86} & 1.97 &  2.66 & 6.85 &  6.86 & $\dagger$ \\
Moderate (Uniform)     & 3.05 &          3.19 & 3.29 &  6.50 & 6.98 & 18.08 & \checkmark \\
Diffuse (Dir~2.0)      & 2.92 &          3.06 & 3.17 &  6.13 & 6.98 & 16.97 & \checkmark \\
Near-Uniform (Dir~10.0)& 3.12 &          3.27 & 3.35 &  6.72 & 6.90 & 20.10 & \checkmark \\
Power-Law ($U^3$)      & 2.23 &          2.52 & 2.43 &  4.08 & 7.09 & 11.77 & $\ddagger$ \\
\bottomrule
\end{tabular}
\vspace{4pt}

\noindent\footnotesize
$\dagger$~Concentrated: M4* (1.86) $<$ M2 (1.89)---M4* \emph{surpasses} the unconstrained baseline.\\
$\ddagger$~Power-law: M4 (2.43) $<$ M4* (2.52)---the quadratic penalty over-penalizes rare distant high-prior sites.\\
Raw data: \texttt{experiments/data/exp7\_prior\_sensitivity\_summary.csv}.
\end{table}

Figure~\ref{fig:prior-sensitivity} and Table~\ref{tab:prior-sensitivity}
report the results. The full ordering holds for three of the five prior
shapes; two qualitatively different exceptions emerge that are themselves
informative.

\emph{Concentrated prior: M4* surpasses the unconstrained baseline.}
Under the Dirichlet(0.2) prior, probability mass concentrates so strongly
on one or two sites that the auction's first anchor choice is almost
certainly the dominant site; the robot travels directly to it and finds
the target with near-minimal FCR regardless of energy budget.
In this regime M4* (mean FCR $= 1.86$) edges out M2 ($1.89$), confirming
the pattern observed in the energy sweep (Section~\ref{sec:energy}) and
the large-fleet regime of the robot sweep: when the search problem is
effectively reduced to a single high-confidence site, the energy budget
is no longer the binding constraint, and the quadratic penalty guides the
anchor placement as accurately as the unconstrained node-to-node auction.
This result extends the energy and fleet-size crossovers to a third
mechanism---prior concentration---under which M4* can surpass M2.

\emph{Power-law prior: M4 edge over M4*.}
The power-law prior ($U(0,1)^3$ normalized) generates a few sites with
very high probability and many with very low probability.
Under this distribution the optimal strategy is to travel to the dominant
high-prior site even if it is geometrically far from base, because the
probability reward is large enough to compensate.
The quadratic penalty of M4* discounts this distant site more than M4's
linear penalty does, causing M4*'s anchor selection to be more
conservative than is warranted by the extreme prior skew.
M4 (mean FCR $= 2.43$) therefore outperforms M4* ($2.52$) by $3.7\%$
on power-law priors.
This exception is not a failure of the cost-foreclosure argument but
rather its boundary: the argument assumes that the anchor site is drawn
from a distribution that spreads probability over the arena without
extreme heavy tails.
In the rare regime of extreme power-law skew, the linear penalty $p/d$
is more appropriate than the quadratic $p/d^2$.
A practical deployment heuristic follows: use M4* when prior distributions
are moderate to concentrated, and revert to M4 when the operational prior
is known to have power-law character.

These two exceptions are theoretically coherent and enrich rather than
undermine the primary findings: the $p/d^2$ bid is optimal over the broad
class of moderate, diffuse, and near-uniform priors that characterize most
real deployments, and it achieves its most striking advantage precisely
when a concentrated prior makes the search problem hard for the
unconstrained baseline.
